\SourceLeaf{362}
\noindent Soll dieser für $t'=0$ gelten, so muss
\[
aC=bD
\]
sein; wodurch
\[
n'=\text{Mod.}\,\frac{a}{D}=\text{Mod.}\,\frac{b}{C}
\]
wird. Soll dieser zugleich $=1$ werden, so kann man, wenn $\varepsilon$ eine beliebige
Grösse, deren Modulus $=1$, ist, setzen
\[
\varepsilon{}C=b,\qquad \varepsilon{}D=a
\]
und
\[
c=\varepsilon{}B,\qquad d=\varepsilon{}A.
\]
Schreibt man anstatt $a,b,c,d$: $\sqrt{\varepsilon}.a$,
$\sqrt{\varepsilon}.b$, $\sqrt{\varepsilon}.c$, $\sqrt{\varepsilon}.d$,
wodurch die Allgemeinheit nicht vermindert wird, so wird schlechthin
\[
c=B,\qquad d=A
\]
oder die Transformationsformel
\[
t'=\frac{at-b}{Bt+A},\qquad
t=\frac{At'+b}{-Bt'+a}.
\]

5. Die einfachste Transformationsformel ist
\[
t'=\frac{t-a}{At+1}.
\]

6. Zwischen zweien Punkten $a,b$ liegt in der Mitte der Punkt
\[
\frac{\BareRoot{(1+aB)(1+bB)}-\BareRoot{(1+aA)(1+bA)}}
{-A\BareRoot{(1+aB)(1+bB)}+B\BareRoot{(1+aA)(1+bA)}}
\]
oder, wenn die den $a,b$ entgegengesetzten Punkte $-\dfrac{1}{A},-\dfrac{1}{B}$
mit $a',b'$ bezeichnet werden,
\[
\frac{a'\BareRoot{(a-b')(b-b')}-b'\BareRoot{(a-a')(b-a')}}
{\BareRoot{(a-b')(b-b')}-\BareRoot{(a-a')(b-a')}}.
\]
